\section{\ledger{}: Extract, Consolidate, Evict by Tier}
\label{sec:method}

If a summariser keeps what was repeated, the fix must take the decision of what to keep away from the summariser.
\ledger{} splits memory maintenance into three steps, of which only the first uses the model.

\paragraph{Extract.}
After session $t$ the model sees \emph{only} the transcript of session $t$, never the memory, and lists the claims worth remembering, one per line, as \texttt{TYPE | SPEAKER | claim} in at most 15 words, with \texttt{TYPE} one of \textsc{warning}, \textsc{condition}, \textsc{dissent}, \textsc{decision}, \textsc{status} or \textsc{other}.
Because the extractor never sees earlier notes it cannot rewrite, merge away or re-confirm them; a claim stated once is extracted once and then left alone.
Lines that do not parse are discarded and counted.

\paragraph{Consolidate.}
Each extracted claim is matched in code against the current items.
It merges into the most similar item if their content-word overlap coefficient is at least $0.4$ and they do not mention different numbers; otherwise it is appended verbatim as a new item.
A merge adds session $t$ to the item's provenance and keeps the original wording.
The number rule keeps successive status updates (``closed 12 tickets'', ``closed 15 tickets'') as separate items so that stale ones can age out.
The threshold was calibrated on the banks: it merges 80--91\% of paraphrase pairs of the same claim and 1 of ${\sim}9.6$k pairs of different claims.

\paragraph{Evict by tier.}
If the rendered items exceed $B$ words, items are removed one at a time in order of the key
$(\mathrm{tier}(\text{type}),\ \text{last mention},\ \text{hash})$, lowest first, where warnings, conditions and dissent form the top tier, decisions the middle tier, and status and other items the bottom tier.
Within a tier the least recently mentioned item goes first; a content hash breaks ties so that the result does not depend on output order.
Support (number of mentioning sessions) is deliberately \emph{not} part of the key: eviction by support is exactly the frequency bias we are trying to remove, and we test it as an ablation.
Decisions are placed below the top tier because they restate the plan that the majority thread already carries; in a flat protected tier they crowd out the rarer warnings when the budget is tight.
The memory is rendered as one bullet per item.

\paragraph{Ablations and baselines.}
All strategies get the same hard word budget and the same reader.
\begin{itemize}
\item \textbf{rolling}: generic rolling summary (\S\ref{sec:benchmark}).
\item \textbf{rolling-imp}: the rolling prompt plus an instruction to preserve every warning, condition, requirement and dissent even if mentioned once, and to drop routine progress first: a deliberately strong prompt-only fix.
\item \textbf{typed}: rewrite under fixed headings (constraints \& warnings, disagreements, decisions, status, other), our re-implementation of type-aware compaction \citep{zerhoudi2026compaction}.
\item \textbf{ace}: untyped notes extracted from the new session, recency eviction, in the spirit of delta updates \citep{zhang2026ace}.
\item \textbf{ledger-supp}: \ledger{} extraction with eviction by support.
\item \textbf{ledger-protected}: one flat protected tier (warnings, conditions, dissent and decisions).
\item \textbf{ledger-prov}: as ledger-protected, rendering provenance tags \texttt{[TYPE; speaker; sessions]}, which cost budget.
\end{itemize}
The extraction output depends only on the session, not on budget or policy, so all itemised variants consolidate exactly the same extracted claims; differences between them are due to consolidation and eviction alone.
\todo{Note asymmetry: itemised strategies are not told B; summary strategies are. Report format-compliance stats (bad lines, empty memories).}
